\section{Integral subspaces and the passage to period pairings}
\label{ip-section}

The local divisor argument applies a valuation bound to every integral
test vector after extending the coefficient field. This passage requires
an equality of integral submodules, not merely an equality of their
complex spans. We prove that equality and its valuation consequence.
The actual lower bounds for the theta periods remain separate inputs;
the following argument specifies exactly what their extension requires.

\subsection{Intersection with a rational subspace after extension}

Let $K_0\subseteq K$ be number fields with rings of integers $A$ and
$B$. Fix an embedding $\sigma:K\hookrightarrow\mathbb C$.
Let $V_0$ be a finite-dimensional $K_0$-vector space, let $M$ be a
full $A$-lattice in $V_0$, and identify
$V_0\otimes_{K_0,\sigma}\mathbb C$ with its specified complex
realization $V$. In particular a $K_0$-basis of $V_0$ is linearly
independent over $\mathbb C$ in this realization. Write $BM$ for
the $B$-span of $M$ in $V$.

\begin{lemma}\label{ip-intersection}
Let $e:V\to V$ be a complex-linear idempotent. Suppose there is
$0\ne a\in A$ such that $ae(M)\subseteq M$. Then
\begin{equation}\label{ip-cut}
                 BM\cap eV=B(M\cap eV).
\end{equation}
This equality holds at ramified primes as well as at unramified primes.
\end{lemma}
\begin{proof}
Since $M$ spans $V_0$, the denominator condition implies that $e$
preserves $V_0$. Set $q=a(1-e)$. It is an $A$-linear endomorphism
of $M$, and its complex-linear extension has kernel $eV$: the
conditions $e^2=e$ and $a\ne0$ show that $qv=0$ is equivalent to
$ev=v$. Thus $\ker(q|_M)=M\cap eV$.

We give the integral base-change argument explicitly. At a nonzero
prime $\mathfrak p$ of $A$, the module $M_{\mathfrak p}$ is free
of finite rank over the discrete valuation ring $A_{\mathfrak p}$.
Choose a basis. Smith normal form, with independent changes of bases
in the source and target, makes the matrix of $q$ diagonal, with
nonzero entries followed by zero entries. Its kernel is the span of
the basis vectors in those zero positions.
The ring $(A\setminus\mathfrak p)^{-1}B$ is a subring of $K$.
Multiplication by a nonzero element of $A_{\mathfrak p}$ remains
injective on it. Therefore the kernel of the same diagonal matrix
after tensoring is still exactly the span of those zero positions.
It follows that the natural map
\[
 B\otimes_A\ker(q|_M)\longrightarrow
          \ker\bigl(1\otimes q:B\otimes_A M\to B\otimes_A M\bigr)
\]
is an isomorphism after localization at every $\mathfrak p$.
It is an isomorphism globally: a nonzero element of its kernel or
cokernel would remain nonzero at a maximal ideal containing its
annihilator, contradicting the local calculation.

The map $B\otimes_A M\to V$ is injective as well. Locally over
$A_{\mathfrak p}$ its chosen basis of $M_{\mathfrak p}$ is a
$K_0$-basis and hence a complex-linearly independent tuple in $V$.
The localized evaluation map is consequently injective; the same
local criterion proves injectivity globally. Its image is $BM$.
Under this identification the last displayed kernel is $BM\cap eV$,
and the image of $B\otimes_A\ker(q|_M)$ is $B(M\cap eV)$.
This proves \eqref{ip-cut}. No assumption about an integral basis
of $B$ over $A$, or about prime splitting, has been made.
\end{proof}

The same proof applies to a finite matrix of maps with a common
denominator: after clearing that denominator, Smith normal form
computes the common kernel locally. This observation is useful if
the desired subspace is specified by linear equations instead of
by an idempotent.

\begin{example}[The subspace must be defined over the original field]
Take $K_0=\mathbb Q$, $K=\mathbb Q(\sqrt2)$,
$M=\mathbb Z^2\subset\mathbb C^2$, and
$e(x,y)=(x,\sqrt2x)$. Then $e^2=e$ and
$eV=\mathbb C(1,\sqrt2)$, but $M\cap eV=0$.
Indeed an integer pair on that line must have both coordinates zero.
On the other hand $(1,\sqrt2)\in BM\cap eV$, so
\eqref{ip-cut} fails. There is no nonzero integer $a$ for which
$ae(M)\subseteq M$. This is a counterexample to omitting the
rationality condition, not to Lemma~\ref{ip-intersection}.
\end{example}

The faithful complex realization is a separate condition. An arbitrary
torsion-free module inside a complex vector space need not satisfy it:
the map $\mathbb C\otimes_{\mathbb Z}(\mathbb Z+\mathbb Z\sqrt2)
\to\mathbb C$ has the nonzero kernel element
$1\otimes\sqrt2-\sqrt2\otimes1$. The source explicitly requires
injectivity of complexification in the hypothesis called an integral
form. Its subsequent intersection argument does not deduce this
injectivity from torsion-freeness alone.

\subsection{The integral denominator of the projectors}

Here is the application of the preceding lemma to the finite-group
subspaces. Let a finite group $G$ act on $V_0$, preserving $M$.
Assume that $K_0$ contains the values of the irreducible characters
in a chosen set $\mathcal X$ and of a one-dimensional character
$\psi:P\to K_0^\times$ of a subgroup $P\subseteq G$.
The two projectors are
\[
 e_{\mathcal X}
   =\frac1{|G|}\sum_{g\in G}\sum_{\chi\in\mathcal X}
                 \chi(1)\chi(g^{-1})\rho(g),
 \qquad
 e_\psi=\frac1{|P|}\sum_{p\in P}\psi(p)^{-1}\rho(p).
\]
The finite-group argument in Section~\ref{div:section} proves that
they are commuting idempotents: the first is central, and the second
projects to the $\psi$-subspace. Character values and roots of unity
are algebraic integers. Consequently, for $e=e_{\mathcal X}e_\psi$,
\begin{equation}\label{ip-denominator}
                     |G|\,|P|\,e(M)\subseteq M.
\end{equation}
Indeed, multiplying the two displayed sums by $|G||P|$ leaves
a finite sum of integral scalar multiples of the operators $\rho(gp)$,
each preserving $M$. Lemma~\ref{ip-intersection} therefore gives
the exact integral intersection with this simultaneous isotypic and
character subspace after extending the coefficient field.
The denominator in \eqref{ip-denominator} is used to make an
integral endomorphism. It is not an extra factor in every test
vector and is not an additional loss in the following valuation bound.

\subsection{Extension of a valuation bound}

For a nonzero prime $\mathfrak q$ of $B$, let
$v_{\mathfrak q}:K^\times\to\mathbb Z$ be the additive valuation
with a uniformizer of value one, and put $v_{\mathfrak q}(0)=+\infty$.
Thus $v_{\mathfrak q}(b)\ge0$ for $b\in B$ and
$v_{\mathfrak q}(x+y)\ge\min(v_{\mathfrak q}(x),v_{\mathfrak q}(y))$.

\begin{proposition}\label{ip-valuation}
Retain the notation of Lemma~\ref{ip-intersection}. Let $U$ be a
finite-dimensional $K$-vector space, let $\mathcal S\subseteq U$,
and let
\[
 \Phi:(U\otimes_{K,\sigma}\mathbb C)\times V\longrightarrow\mathbb C
\]
be complex-bilinear. Fix $k\in\mathbb Z$. Suppose that for every
$s\in\mathcal S$ and $m\in M\cap eV$ there is $a_{s,m}\in K$
with
\[
       \Phi(s,m)=\sigma(a_{s,m}),\qquad
                         v_{\mathfrak q}(a_{s,m})\ge k.
\]
Then for every $x\in B\mathcal S$ and $t\in BM\cap eV$ there is
a unique $a_{x,t}\in K$ with
\[
       \Phi(x,t)=\sigma(a_{x,t}),\qquad
                         v_{\mathfrak q}(a_{x,t})\ge k.
\]
The same assertion holds for the $B_{\mathfrak q}$-spans of both
sets, with the evident localization of the test lattice.
\end{proposition}
\begin{proof}
Write $x=\sum_i b_is_i$ with $b_i\in B$, and use
\eqref{ip-cut} to write $t=\sum_j c_jm_j$ with $c_j\in B$ and
$m_j\in M\cap eV$. Bilinearity gives
\[
       \Phi(x,t)=\sigma\left(\sum_{i,j}b_ic_ja_{s_i,m_j}\right).
\]
Each summand on the right has valuation at least $k$, since
both coefficients have nonnegative valuation. Their sum has
valuation at least $k$ as well. Injectivity of $\sigma$ proves
uniqueness, so the value does not depend on the chosen expansions.
For localization the same proof applies with coefficients in
$B_{\mathfrak q}$, whose valuations are still nonnegative.
\end{proof}

\begin{corollary}\label{ip-scaling}
Suppose a $K$-valued bilinear period pairing takes
$M_{\mathfrak q}\times T_{\mathfrak q}$ into
$\mathfrak q^{-b}B_{\mathfrak q}$, where $b\in\mathbb Z$.
If $x\in\mathfrak q^aM_{\mathfrak q}$ with $a\in\mathbb Z$,
then its pairing against every $t\in T_{\mathfrak q}$ has
valuation at least $a-b$.
\end{corollary}
\begin{proof}
Choose a uniformizer $\pi$ of $B_{\mathfrak q}$ and write
$x=\pi^a y$ with $y\in M_{\mathfrak q}$. The pairing equals
$\pi^a\Phi(y,t)$, so its valuation is at least $a-b$.
This argument uses only a local generator; it makes no assertion
that the global prime ideal is principal.
\end{proof}

These results justify the transition from integral period tests
to all vectors in the final coefficient lattice.  In the application
of~\cite{alpo-uniform-source}, the projection is the level operator,
its integral denominator is the finite group order, and the cut
condition is \eqref{ip-cut}.

What remains is mathematical content, not another scalar-extension
lemma: one must prove the initial valuation bound for the actual
theta periods and prove the asserted containment of their actual
source vectors in the required local lattices. The submitted
construction distinguishes the level prime, the two auxiliary
primes, and the other primes, with separate estimates in each case.
Proposition~\ref{ip-valuation} does not supply any of those estimates.


